\section{Operations with Landau Symbols}
\label{sec:landau}

The
\index{Landau symbols}%
\index{Landau!symbols of}%
\index{small $o$}%
\index{small o}%
\index{big $O$}%
\index{big O}%
Landau notations, small $o$ and big $O$, are extremely convenient for handling
asymptotic estimates. However, one must be very careful with how these notations
are used, as otherwise, one risks making gross errors. Some counterintuitive
results are, for example:
\begin{enumerate}
\item $o(x) - o(x) = o(x)$ (and not $0$),
\item $o(x^2) = o(x)$ but not $o(x) = o(x^2)$.
\end{enumerate}

The second example, in particular, tells us that the equality symbol is not
actually used appropriately in this context, as it does not possess the
symmetric property. In this section, we propose a formally precise definition of
the small $o$ (and big $O$) object and a rigorous way to manipulate and compare
it.

\begin{definition}[small $o$, big $O$]
Let $A\subset \RR$, $x_0$ be an accumulation point of $A$, and let $g\colon A
\to \RR$ be a positive function on $A$. We then define the sets of
functions\mynote{%
recall that $A^B$ denotes the set of functions $f\colon B \to A$}:
\begin{align*}
  o(g) &= \ENCLOSE{f\in \RR^A \colon \lim_{x\to x_0}\frac{f(x)}{g(x)} = 0};\\
    O(g) &= \ENCLOSE{f\in \RR^A \colon \limsup_{x\to x_0}\abs{\frac{f(x)}{g(x)}}
  < + \infty}\\
              &= \ENCLOSE{f\in \RR^A \colon \exists U\in \U_{x_0}, C>0 \colon
       \forall x \in U\colon \abs{\frac{f(x)}{g(x)}}\le C}.
\end{align*}
\end{definition}

Expressions such as:
\[
  \sin x - x = o(x^2), \qquad \text{for $x\to 0$}
\]
should thus be interpreted as
\[
  f \in o(g)
\]
where $f(x) = \sin x -x $ and $g(x) = x^2$.
In this section, we will write:
\begin{equation}\label{eq:39523}
  \sin x - x \in o(x^2)
\end{equation}
to emphasize this interpretation (and to highlight the fact that the relation is
not symmetric at all).

We can now proceed to interpret operations between sets. In general, if $A$ and
$B$ are sets of objects on which an operation $*$ is defined (which could be
addition, multiplication, division...), we define:
\[
  A * B = \ENCLOSE{a*b \colon a \in A, b\in B}.
\]
Similarly, if $A$ is a set on whose elements an operation $*$ with an object $b$
is defined, we will define:
\[
  A * b = \ENCLOSE{a*b\colon a \in A}, \qquad
  b * A = \ENCLOSE{b*a\colon a \in A}.
\]
For example, the expression
\[
  \sin x = x + o(x^2)
\]
should formally be understood as
\[
  \sin x  \in x + o(x^2)
\]
where the set $x+ o(x^2)$ is the set of all functions
$f$ that can be written in the form $f(x) = x+h(x)$ with $h\in o(x^2)$, i.e.,
$h$ such that $h(x)/x^2 \to 0$. It follows that this expression is equivalent to
\eqref{eq:39523} since if $\sin x = x + h(x)$ with $h\in o(x^2)$, it means that
$\sin x- x \in o(x^2)$.

We can then state the algebraic rules for manipulating Landau symbols.

\begin{theorem}[Operations with Landau Symbols]
Let $A\subset \RR$, $x_0\in [-\infty, +\infty]$ be an accumulation point for
$A$. Let $f,g \colon A \to \RR$ be positive functions, $c\in \RR$, $c\neq 0$,
$n\in \NN$, $n\neq 0$. Then, for $x\to x_0$, we have:
\begin{align*}
0 \in o(f) &\subset O(f) & f&\in O(f)& \\
c \cdot o(f) &= o(f) & c \cdot O(f) &= O(f)\\
o(f)+o(f) &= o(f) & O(f)+O(f) &= O(f)\\
o(f)-o(f) &= o(f) & O(f)-O(f) &= O(f)\\
o(f\cdot g) &= f \cdot o(g) & O(f\cdot g) &= f\cdot O(g)\\
o(g/f) &= o(g) / f & O(g/f) &= O(g) / f\\
o(f)\cdot o(g) &\subset o(f\cdot g) & O(f)\cdot O(g) &\subset O(f\cdot g) \\
o(f)\cdot O(g) & \subset o(f\cdot g)&  & \\
(o(f))^n &\subset o(f^n) & (O(f))^n &\subset O(f^n)\\
o(o(f)) &\subset o(f) & O(O(f)) &\subset O(f)\\
o(O(f)) &\subset o(f) & O(o(f)) &\subset o(f)
.
\end{align*}
It is observed, in particular, that $o(f)$ and $O(f)$ are vector subspaces of
$\RR^A$.

It is also possible to apply variable changes.
If $f\in o(g)$ for $x\to x_0$ and $h(t) \to x_0$ for $t\to t_0$, then
\[
  f \circ h \in o(g\circ h) \qquad \text{for $t\to t_0$}.
\]
\end{theorem}
%
\begin{proof}
Since $0/f = 0$, it is obvious that $0\in o(f)$.
The inclusion $o(f) \subset O(f)$ follows from the fact that if $h\in o(f)$, it
means that $h/f\to 0$ for $x\to x_0$. But then there exists a neighborhood of
$x_0$ in which $\abs{h/f}<1$, and thus $h\in O(f)$. Obviously, $f\in O(f)$ since
$f/f=1$ is a bounded function.

If $h\in c \cdot o(f)$, it means that $h = c \cdot (h/c)$ and $h/c \in o(f)$,
i.e., $(h/c)/f \to 0$. But this is equivalent to $h/f\to 0$ since $c$ is a
non-zero constant. The case of big $O$ is handled similarly.

The set $o(f)+o(f)$ consists of functions of the form $h+k$ with $h,k \in o(f)$.
But we have
\[
  \frac{h+k}{f} = \frac{h}{f} + \frac{k}{f} \to 0 + 0 = 0.
\]
Thus $o(f)+o(f) \subset o(f)$. Conversely, observe that if $h\in o(f)$, it can
be written as $h = h/2 + h/2$, and from what has been seen before, we know that
$h/2 \in o(f)$. Thus $o(f) \subset o(f) + o(f)$.
A similar reasoning can be applied to $O(f)$: the sum of two locally bounded
functions is also locally bounded.

Observe that $o(f)-o(f) = o(f) + (-1)\cdot o(f) = o(f) + o(f) = o(f)$
from what has already been seen. The same holds for $O(f)-O(f)$.

If $h\in o(fg)$, it means that $h/(fg)\to 0$. But then $h = f\cdot (h/f)$ with
$h/f \in o(g)$ since $(h/f)/g = h/(fg)\to 0$. Thus $h \in fo(g)$, and
consequently $o(fg)\subset f \cdot o(g)$. Conversely, if $h \in f\cdot o(g)$, it
means that $h = f \cdot(h/f)$ with $h/f \in o(g)$, i.e., $(h/f)/g \to 0$. Thus
$h/(fg) \to 0$, which means $h\in o(fg)$.
A similar reasoning can be applied to big $O$.

The results for division are reduced to those of multiplication by observing
that division by $f$ is equivalent to multiplication by $1/f$.

If $h\in o(f)$ and $k\in o(g)$, we want to show that $hk\in o(fg)$.
But this is obvious since $(hk)/(fg) = (h/f)\cdot(k/g) \to 0\cdot 0 = 0$. We
have shown that $o(f)o(g)\subset o(fg)$. A similar result holds in the other two
cases $O(f)O(g)\subset O(fg)$ and $o(f)O(g)\subset o(fg)$, observing that the
product of two bounded functions is bounded, and the product of a bounded
function with an infinitesimal is infinitesimal.

Regarding the power $(o(f))^n$, observe that:
\begin{align*}
  (o(f))^n &= \ENCLOSE{h^n \colon h\in o(f)}
      \subset \ENCLOSE{h_1 \cdots h_n \colon h_1, \dots, h_n \in o(f)}\\
      &= o(f) \cdots o(f) \subset o(f^n).
\end{align*}

If $h\in o(o(f))$, it means that there exists $k\in o(f)$ and $h\in o(k)$. Thus
$k/f\to 0$ and $h/k\to 0$. But then $h/f = (h/k)\cdot (k/f) \to 0\cdot 0 = 0$,
thus $h \in o(f)$. We have therefore shown that $o(o(f)) \subset o(f)$. A
similar demonstration can be made for big $O$. The inclusions $o(O(f))\subset
o(f)$ and $O(o(f))\subset o(f)$ are demonstrated similarly, observing that the
product of a bounded function with an infinitesimal is infinitesimal.

Regarding the change of variable, we need to verify that
\[
  \frac{f(h(t))}{g(h(t))} \to 0
  \qquad\text{for $t\to t_0$}
\]
if $f\in o(g)$ and if $h(t)\to x_0$ for $t \to t_0$. But this is nothing other
than the change of variable $x=h(t)$ in the limit $f(x)/g(x)\to 0$ for $x\to
x_0$.
\end{proof}

\begin{example}
Knowing that for $x\to 0$, we have
\[
\sin x \in x + o(x^2), \qquad \cos x \in 1 - \frac {x^2}{2} + o(x^3)
\]
and recalling that $x^3 \in o (x^2)$, we can deduce, using the properties of the
previous theorem:
\begin{align*}
2\cos x  - \sin x
&\in 2 - x^2 + 2o(x^3) - x - o(x^2)\\
&= 2- x - x^2 + o(x^3) + o(x^2)\\
&= 2- x - x^2 + o(o(x^2)) + o(x^2)\\
 &\subset 2 - x - x^2 + o(x^2) + o(x^2)\\
 &= 2-x-x^2 + o(x^2).
\end{align*}
Usually, all these steps are implied, the symbol $=$ is used instead of $\in$
and $\subset$, and often only one $o(\cdot)$ is preferred at the end of each
expression.
\end{example}

\begin{example}
With the same assumptions as the previous example, we could write:
\begin{align*}
(\sin^2 x)\cdot(2-2\cos x)^3
&\in (x+o(x^2))^2 \cdot (2 - (2- x^2 - 2o(x^3))^3 \\
&= (x+o(x^2))^2 \cdot (x^2+o(x^3))^3 \\
&= (x^2 + 2 x o(x^2) + (o(x^2))^2) \\
&\quad \cdot (x^6 + 3 x^4 o(x^3) + 3 x^2(o(x^3))^2 + (o(x^3))^3)\\
&\subset (x^2 + o(x^3) + o(x^4))\cdot(x^6 + o(x^7) + o(x^8) + o(x^9)) \\
&= (x^2 + o(x^3)) \cdot(x^6+o(x^7)) \\
&= x^8 + x^2o(x^7) + o(x^3) x^6 + o(x^3)o(x^7)\\
&= x^8 + o(x^9) + o(x^9) + o(x^{10})
= x^8 + o(x^9).
\end{align*}
\end{example}

\begin{example}
We can also perform an exercise with the change of variable.
Observing that $y=1-\cos x\to 0$ for $x\to 0$
and knowing that $\sin y = y + o(y^2)$ for $y\to 0$
and that $2-2\cos x = x^2 + o(x^2)$ for $x\to 0$,
we can write
that for $x\to 0$:
\begin{align*}
\sin(2-2\cos x)
& \in (2-2\cos x) + o((2-2\cos x)^2) \\
& \subset x^2 + o(x^2) + o((x^2+o(x^2))^2) \\
& \subset x^2 + o(x^2) + o(O(x^4)) \\
& \subset x^2 + o(x^2) + o(x^4) \\
& = x^2 + o(x^2).
\end{align*}
\end{example}

\begin{exercise}
Prove (out of simple curiosity) that the inverse inclusions to those stated in
the theorem also hold:
\begin{gather*}
o(f\cdot g) \subset o(f) \cdot o(g), \qquad
o(f\cdot g) \subset O(f) \cdot o(g), \\
o(f\cdot g) \subset o(f) \cdot O(g), \\
o(f) \subset o(o(f)),  \qquad
O(f) \subset O(O(f)).
\end{gather*}
Observe instead that $(o(f))^2 \neq o(f)\cdot o(f)$.
\end{exercise}

The Taylor formula can be very useful for determining the nature of a series, as
in the following.
\begin{exercise}
Determine the values of $\alpha \in \RR$ for which the series
\[
  \sum_{k=1}^{+\infty}(-1)^k \enclose{\frac 1 k - \sin \frac 1 k}^\alpha
\]
converges absolutely and those for which it converges.
\end{exercise}
\begin{proof}[Solution.]
Let $f(x) = x - \sin x$,
through Taylor expansion, we know that,
for $x\to 0$:
\begin{align*}
   f^\alpha(x) &= \enclose{x - \sin x}^\alpha
   =  \enclose{\frac{x^3}{6} + o(x^3)}^\alpha
   = \frac{x^{3\alpha}}{6^\alpha}\enclose{1+o(1)}^\alpha\\
   & = \frac{x^{3\alpha}}{6^\alpha}\enclose{1 + \alpha \cdot o(1) + o(o(1))}
   = \frac{x^{3\alpha}}{6^\alpha}\enclose{1 + o(1)}.
\end{align*}
Since for $k\to +\infty$, we have $x=1/k \to 0$, we can
write:
\[
  a_k = \enclose{f(1/k)}^\alpha = \enclose{\frac 1 k - \sin \frac 1 k}^\alpha
  = \frac{1}{6^\alpha \cdot k^{3\alpha}} \enclose{1+o(1)}
  \sim \frac{1}{6^\alpha \cdot k^{3\alpha}}.
\]
Observe that the given series is $\sum (-1)^k a_k$ and that $a_k>0$. Thus, by
the asymptotic comparison test, if $3\alpha > 1$, i.e., if $\alpha >1/3$,
the given series converges absolutely, whereas if $\alpha \le 1/3$, the series
does not converge absolutely (but it might converge).
Also note that if $\alpha \le 0$, the term $a_k$ is not even infinitesimal, and
thus the series cannot converge.

For $\alpha \in (0,1/3]$, we can try to use the Leibniz criterion:
the sequence $a_k$ is infinitesimal, and we need to verify that it is also
eventually decreasing.
Having set $a_k = \enclose{f(1/k)}^\alpha$, it will be sufficient to show that
the function $f(x)$ is increasing in a right neighborhood of $0$.
By the properties of the Taylor polynomial, we know that the Taylor polynomial
of order $2$ of $f'(x)$ is equal to the derivative of the Taylor polynomial of
$f(x)$ of order $3$. Thus, we can immediately assert (but it would have been
equally quick to calculate the derivative and then expand it) that
\[
  f'(x)
  = \enclose{\frac{x^3}{6}}'  + o(x^2)
  = \frac{x^2}{2} + o(x^2) = x^2 \enclose{\frac 1 2 + o(1)}
\]
thus, by the permanence of the sign, there must exist a neighborhood of $0$ in
which
$1/2 + o(1)$ is positive, and therefore in such a neighborhood, it results that
$f'(x)\ge 0$
i.e., $f$ is increasing.
Thus, $f^\alpha$ is also increasing, and for sufficiently large $k$,
$a_k$ is decreasing. We can then apply the Leibniz criterion and
conclude that the series is convergent for every $\alpha >0$.
\end{proof}

\begin{exercise}
  Prove that the following series is convergent:
  \[
    \sum_k \sqrt{k} \enclose{k\tg{\frac 1 k} - \cos \frac 1 k}.
  \]
\end{exercise}
\begin{proof}[Solution]
In this exercise, we use the notion of big $O$ (but the exercise could also be
solved using small $o$, as usual).
For $x\to 0$, we have: %
\mynote{%
here we exploit the fact that the third-order Taylor expansion exists 
and using the big $O$ notation, we avoid having to remember the coefficient
$c=\frac 1 3$ that appears in the formula}
\[
  \tg x = x + c x^3 + o(x^3) = x + O(x^3)
\]
and similarly
\[
  \cos x = 1 + O(x^2)
\]
from which
\[
  \frac{\tg x}{x} - \cos x = O(x^2).
\]
Thus, for $k\to +\infty$, we have
\[
  a_k = \sqrt{k} \enclose{k\tan{\frac 1 k} - \cos \frac 1 k}
      = \sqrt{k} \cdot O\enclose{\frac 1 {k^2}} = O\enclose{\frac 1 {k^{\frac 3
   2}}}.
\]
This means that for sufficiently large $k$, there exists a constant $C$
such that
\[
  a_k \le \frac{C}{k^{\frac 3 2}}.
\]
Since we know that the series $\sum \frac{1}{k^p}$ is convergent
if $p>1$, by the comparison test, $\sum a_k$ (which is a series with positive
terms) is also convergent.
\end{proof}